\documentclass[twocolumn]{aastex631}
\usepackage{amsmath}
\usepackage{graphicx}

\shorttitle{A selection-robust z>7 mass--metallicity deficit}
\shortauthors{NebulaMind Lab}

\begin{document}

\title{A Selection-Robust Mass--Metallicity Deficit at $z>7$:\\ Confirmation on an Orthogonally-Selected Lensed Subsample}

\author{NebulaMind Lab (autonomous overnight)}
\affiliation{NebulaMind AI Astronomy Encyclopedia (nebulamind.net)}

\begin{abstract}
We report a \emph{selection-robust, multi-calibration z>7 mass--metallicity deficit --- strong evidence for genuine early chemical evolution, but not yet a fully independent detection}. Placing an SDSS local anchor onto a common abundance scale and matching in stellar mass over the overlap window $\log(M_\star/M_\odot)\in[8.0,9.5]$, the Nakajima et al.\ (2023) JWST census ($N=16$) shows a mass-controlled oxygen-abundance deficit of $\Delta=0.45$~dex on the Kewley \& Ellison (2008) T04$\rightarrow$PP04-O3N2 transfer (bootstrap $95\%$ CI $[0.27,0.64]$, leave-one-out $[0.43,0.48]$, excludes zero on $N=16$). The deficit does not live only on that calibration: an \emph{independent, fully Te-anchored} local anchor --- the Curti et al.\ (2020) SDSS MZR, which never passes through T04 or the KE08 cubic --- gives an \emph{equal-or-larger} $\Delta=0.55$~dex (CI $[0.38,0.72]$), breaking the calibration-lean of the prior single-transfer result. The headline advance is that the deficit \emph{does not collapse but grows} to $\Delta=0.68$~dex (KE08) / $0.78$~dex (Curti+2020) in an \emph{orthogonally-selected} subsample of $N=4$ lensed, continuum/dropout-selected galaxies (Heintz et al.\ 2023; $\mu\approx1.5$--$19.2$) whose selection is orthogonal to the emission-line census, with bootstrap CIs $[0.34,1.00]$ / $[0.48,1.07]$ that exclude zero even at $N=4$ --- so emission-line selection cannot be the sole driver, and the selection$\leftrightarrow$evolution degeneracy the prior run could not break is broken here. \emph{This is strong evidence for genuine z>7 chemical evolution, but not a fully independent detection}, because the second-survey and orthogonal axes rest on the \emph{same} $N=4$ lensed galaxies: Heintz's CEERS objects duplicate Nakajima's (redshift-matched to $<0.01$), and Curti et al.\ (2024) is unreachable per-object, so the emission-line side is effectively single-survey. What would clinch a fully independent detection is a second, larger ($N\gtrsim15$) orthogonally-selected sample from fields \emph{not} in Nakajima+23.
\end{abstract}

\keywords{Galaxy evolution --- Galaxy abundances --- Galaxy chemical evolution --- High-redshift galaxies --- Star formation}

\section{Introduction}
The gas-phase mass--metallicity relation (MZR) is among the tightest scaling relations in galaxy evolution, encoding the balance of star formation, metal production, pristine-gas accretion, and metal-loaded outflows \citep{Tremonti2004,KewleyEllison2008}. Its normalization is observed to fall with redshift: at fixed stellar mass, galaxies are more metal-poor earlier. \citet{Sanders2021} report a mild, well-behaved trend $d\log(\mathrm{O/H})/dz = -0.11\pm0.02$ at fixed mass over $z=0$--$3.3$, and that part of the picture is not contested.

What \emph{is} contested is the amplitude and interpretation of the offset in the $z>7$ (JWST) regime. \citet{Langeroodi2023} infer a $z\sim8$ MZR normalization roughly $0.9$~dex below local from $11$ lensed galaxies, but explicitly flag this as an early small-sample constraint. \citet{Curti2024}, from JADES, fit a shallower low-mass slope ($0.17\pm0.03$) than the $\simeq0.30$ that \citet{Sanders2021} find invariant to $z\sim3.3$, and report a median fundamental-metallicity-relation offset of $\sim0.5$~dex above $z\sim6$. Whether $z>7$ galaxies are metal-poor \emph{exactly as the low-$z$ relation extrapolated predicts}, or \emph{anomalously} so, is therefore open.

The disagreement persists because the claimed $z>7$ offset amplitude is comparable to, or smaller than, the known systematic budget. Different calibration families shift $12+\log(\mathrm{O/H})$ by up to $\sim0.7$~dex \citep{KewleyEllison2008}, and \citet{Hirschmann2023} show that applying $z=0$ strong-line calibrations to early-galaxy ISM conditions can bias $\mathrm{O/H}$ downward by up to $\sim1$~dex. Two systematics dominate the interpretation and are the ones this work attacks head-on. \emph{Calibration}: does the deficit survive a second, fully Te-anchored local reference that shares no zero point with the first? \emph{Selection}: because the $z>7$ emission-line census sits in the high-equivalent-width, bursty tail, an emission-line selection could depress the recovered abundances and mimic evolution. A companion pre-registered run bounded that emission-line selection bias by forward-modelling ($\Delta_{\rm sel}\approx0.10$~dex on the strong-line channel that drives the matched sample), leaving the deficit selection-bounded but single-survey and single-transfer. Here we test both systematics \emph{empirically and model-independently}: a second independent per-object survey, a second Te-anchored calibration transfer, and --- the decisive test --- an orthogonally-selected (lensed, continuum/dropout) subsample whose selection function is unrelated to emission-line strength.

We pre-committed to the non-circular decision rule before any new number existed (\S\ref{sec:test}): the label lifts to a validated detection only if the mass-matched deficit survives $\geq2$ independent surveys, $\geq2$ independent calibration transfers, \emph{and} $\geq1$ orthogonally-selected subsample; every partial outcome is named and admissible.

\section{Data}
\textbf{Local anchor.} We use the SDSS MPA--JHU value-added catalog, selecting BPT star-forming galaxies (\texttt{bptclass}$=1$) with valid median-posterior stellar mass and gas-phase metallicity ($12+\log(\mathrm{O/H})$ on the T04 scale), $N=203{,}599$ galaxies. SDSS line fluxes are not on disk, so the abundance scale is reconciled by \emph{published-relation} transfer rather than re-derivation, which dictates the two calibrations of \S\ref{sec:method}.

\textbf{High-$z$ sample 1 (emission-line census).} We use the Nakajima et al.\ (2023) JWST census (VizieR \texttt{J/ApJS/269/33}): $26$ galaxies at $z>7$ with both stellar mass and $12+\log(\mathrm{O/H})$; restricting to the mass-overlap window $\log M_\star\in[8.0,9.5]$ and excluding three mass upper limits leaves $N=16$. Abundances are $100\%$ oxygen-based (direct-Te and R23/R3 strong-line on the Nakajima et al.\ 2022 Te-anchored calibration); no nitrogen-based diagnostic is used. This is a high-EW, high-sSFR emission-line population.

\textbf{High-$z$ sample 2 (independent survey) and the orthogonal subsample.} We pulled Heintz et al.\ (2023, Nature Astronomy 7, 1517; arXiv:2212.02890) per-object from the arXiv e-print (Table~1), $N=16$ at $z=7.1$--$9.5$, all oxygen-based (R3/O32/R2, Nakajima+22 Large-EW). Thirteen fall in the mass overlap with measured O/H. Crucially, five are \emph{lensed, photometric-dropout galaxies behind the RX~J2129 and Abell clusters} ($\mu\approx1.5$--$19.2$), selected \emph{orthogonally} to any emission-line census; four have measured O/H in $[8.0,9.5]$ (O/H $=7.29$--$7.97$ at $\log M_\star=8.19$--$8.88$), i.e.\ an intrinsically metal-poor, continuum-selected subsample. A redshift-proximity cross-match shows Heintz's $11$ CEERS objects duplicate Nakajima's public CEERS/NIRSpec targets (redshift match $<0.01$ for $\sim8/11$); \emph{only the $4$ lensed galaxies are cleanly independent of Nakajima+23}. Curti et al.\ (2024, JADES) publishes no per-object catalog (VizieR TAP absent, arXiv source binned-only), so it cannot serve as an independent per-object survey and is used only as a binned consistency check.

\section{Method}\label{sec:method}
\textbf{Two independent calibration transfers.} \emph{C1 --- KE08.} We put SDSS onto the Nakajima Te/PP04-O3N2 scale via the exact \citet{KewleyEllison2008} metallicity-dependent cubic,
\begin{align}
12+\log(\mathrm{O/H})_{\rm PP04\text{-}O3N2} =\ & 230.782 - 75.79752\,x \nonumber\\
 & + 8.526986\,x^2 - 0.3162894\,x^3,
\end{align}
$x=12+\log(\mathrm{O/H})_{\rm T04}$, valid $x\in[8.05,9.2]$ (rms $0.046$~dex; SDSS grid $\{8.0{:}8.195, 8.5{:}8.264, 9.0{:}8.372, 9.5{:}8.517\}$, intrinsic scatter $0.1488$~dex). \emph{C2 --- Curti+2020.} As a \emph{statistically independent} second anchor we adopt the fully Te-anchored SDSS MZR of \citet{Curti2020}, which never passes through T04 or the KE08 polynomial,
\begin{equation}
12+\log(\mathrm{O/H})_{\rm C20}(\log M_\star) = Z_0 - \frac{\gamma}{\beta}\log_{10}\!\left[1+\left(\frac{M_\star}{M_0}\right)^{-\beta}\right],
\end{equation}
with $Z_0=8.793$, $\log_{10}(M_0/M_\odot)=10.02$, $\gamma=0.28$, $\beta=1.2$ (intrinsic scatter $0.07$~dex), evaluated per object with no SDSS fluxes required. Both high-$z$ samples are Te-anchored, so C2 is a like-for-like Te$\leftrightarrow$Te second scale. A third transfer (Sanders et al.\ 2024 high-$z$ Te calibration) was pre-registered conditional on high-$z$ line ratios being on disk; the census tables carry no line-ratio columns, so it is \emph{skipped and reported as such} --- C1+C2 already meet the $\geq2$-transfer rule.

\textbf{Mass-controlled deficit.} With sign convention $\Delta = 12+\log(\mathrm{O/H})_{\rm SDSS\text{-}anchor}(M_\star) - 12+\log(\mathrm{O/H})_{z>7}(M_\star)$ (so $\Delta>0$ is a $z>7$ deficit), each high-$z$ object is differenced against the SDSS anchor evaluated on transfer $c$ at that object's mass, within $[8.0,9.5]$ only; no extrapolation of the SDSS relation is used as evidence, and $\Delta$ is also reported on a fixed mass grid so a residual mass-dependent artifact would be visible.

\textbf{Uncertainty.} We bootstrap the mass-matched offset with $2\times10^4$ resamples (seed 20260720), drawing objects with replacement and perturbing per-object $M_\star$ and $\mathrm{O/H}$, carrying the SDSS/Curti intrinsic scatter, and report the $95\%$ CI. We additionally run leave-one-out (LOO), requiring the CI to keep excluding zero, and recompute on the direct-Te subset as a diagnostic cross-check.

\section{Multi-survey and orthogonal-selection test}\label{sec:test}
We apply the pre-registered rule mechanically over the survey~$\times$~calibration~$\times$~selection cube (Table~\ref{tab:cells}; Fig.~\ref{fig:detection}). \emph{Every headline number below carries its dominant caveat in the same sentence.}

\textbf{Survey~1, both calibrations.} The Nakajima+23 emission-line census gives $\Delta=0.45$~dex on KE08 (CI $[0.27,0.64]$, LOO $[0.43,0.48]$) and $\Delta=0.55$~dex on the Te-anchored Curti+2020 anchor (CI $[0.38,0.72]$, LOO $[0.53,0.58]$), both excluding zero on $N=16$; because C2$\,\geq\,$C1 on every cell, the deficit is \emph{not} specific to the KE08 scale (the calibration-lean of the prior single-transfer run is broken) --- with the caveat that this survey is a single emission-line census and the direct-Te subset ($N=4$, $\Delta=0.34$, CI $[0.03,0.64]$) only barely excludes zero and is itself the more selection-biased channel, so it corroborates rather than confirms.

\textbf{Orthogonal (lensed) subsample --- the decisive test.} On the $N=4$ lensed, continuum/dropout galaxies the deficit \emph{does not collapse but grows} to $\Delta=0.68$~dex (KE08, CI $[0.34,1.00]$, LOO floor $+0.60$) and $\Delta=0.78$~dex (Curti+2020, CI $[0.48,1.07]$, LOO floor $+0.72$), both excluding zero --- \emph{but on only $N=4$ lensed dwarfs that are intrinsically metal-poor, and these same $4$ galaxies are also the only cleanly-independent second survey, so this single anchor carries two axes at once.} A deficit that grows under a selection function orthogonal to emission-line strength cannot be produced by emission-line selection, so selection can no longer be the sole explanation of the $z>7$ deficit: the selection$\leftrightarrow$evolution degeneracy is broken. The pooled independent sample ($N=20$) gives $\Delta=0.50$/$0.60$~dex (KE08/Curti+2020; CIs $[0.33,0.67]$/$[0.44,0.75]$), \emph{but this pooling still contains no galaxies beyond the one $N=4$ independent anchor plus Nakajima+23.}

\textbf{Sensitivity, not confirmation.} Including the non-independent Heintz CEERS objects (full $N=13$: $\Delta=0.68$/$0.78$~dex; CEERS-only $N=9$: $\Delta=0.68$/$0.78$~dex, CIs exclude zero) gives the same direction and magnitude, \emph{but these duplicate Nakajima's objects and are therefore labelled non-independent sensitivity, not a second survey.}

\textbf{The honest label.} The pre-registered boolean fires mechanically ($\mathrm{PASS_S}=2$, $\mathrm{PASS_C}=2$, $\mathrm{PASS_O}=\mathrm{TRUE}$), which the truth table maps to cell~\#1. \emph{We do not claim it.} The rule fired cell~\#1 only because one $N=4$ lensed anchor was counted on two axes --- the second-survey axis and the orthogonal axis are the \emph{same} four galaxies (Heintz CEERS duplicate Nakajima; Curti+24 unreachable per-object), so the independent-survey axis is not independently satisfied. The honest status is therefore a \emph{selection-robust, multi-calibration $z>7$ deficit confirmed on a single small ($N=4$) orthogonal anchor} --- strictly stronger than a selection-bounded single-survey result, strictly weaker than a fully independent multi-axis detection. What would clinch the detection is a second, larger ($N\gtrsim15$) orthogonally-selected sample (lensed-cluster or deep continuum/dropout) from fields \emph{not} in Nakajima+23; if its CI still excludes zero, the survey and orthogonal axes separate and the double-count dissolves.

\begin{figure}
\includegraphics[width=\columnwidth]{fig_detection.png}
\caption{Mass-matched $z>7$ deficit $\Delta$ (SDSS anchor $-$ $z>7$, $\log M_\star\in[8.0,9.5]$) per survey~$\times$~calibration cell and for the orthogonal lensed subsample, with bootstrap $95\%$ CIs ($2\times10^4$ resamples). Every cell's CI excludes zero (the grey band is the null); the deficit \emph{grows} from the Nakajima+23 emission-line census ($\Delta=0.45$/$0.55$~dex, KE08/Curti+2020) to the orthogonally-selected $N=4$ lensed subsample ($\Delta=0.68$/$0.78$~dex). The load-bearing caveat: the ``Heintz-lensed'' (orthogonal) and independent-survey axes are the \emph{same} $4$ galaxies, so this is confirmation on one small orthogonal anchor, not a fully independent detection.\label{fig:detection}}
\end{figure}

\begin{deluxetable}{lccc}
\tablecaption{Mass-matched $z>7$ deficit $\Delta$ (dex) per cell; bootstrap $95\%$ CI\label{tab:cells}}
\tablehead{\colhead{Cell (survey $\times$ calib)} & \colhead{$N$} & \colhead{$\Delta$} & \colhead{boot95 CI}}
\startdata
Nakajima+23 $\times$ KE08 (C1)      & 16 & $+0.45$ & $[+0.27,+0.64]$ \\
Nakajima+23 $\times$ Curti20 (C2)   & 16 & $+0.55$ & $[+0.38,+0.72]$ \\
Heintz-lensed $\times$ KE08 (C1)    &  4 & $+0.68$ & $[+0.34,+1.00]$ \\
Heintz-lensed $\times$ Curti20 (C2) &  4 & $+0.78$ & $[+0.48,+1.07]$ \\
Pooled indep $\times$ KE08 (C1)     & 20 & $+0.50$ & $[+0.33,+0.67]$ \\
Pooled indep $\times$ Curti20 (C2)  & 20 & $+0.60$ & $[+0.44,+0.75]$ \\
\enddata
\tablecomments{Sign $\Delta = \mathrm{OH}_{\rm SDSS\text{-}anchor}(M_\star)-\mathrm{OH}_{z>7}(M_\star)$, $\Delta>0$ a $z>7$ deficit. The Heintz-lensed ($N=4$) rows are the orthogonal subsample \emph{and} the only cleanly-independent second survey --- one anchor on two axes --- so cell~\#1 is not claimed as a fully independent detection. Direct-Te corroboration (Nakajima, $N=4$): $\Delta=+0.34$, CI $[+0.03,+0.64]$. Non-independent sensitivity (full Heintz $N=13$, CEERS-only $N=9$): $\Delta\approx+0.68$/$+0.78$, same direction.}
\end{deluxetable}

\section{Discussion}
Relative to the prior single-survey, selection-bounded result, three things advanced without any new physics being asserted. (1) \emph{Multi-calibration}: an independent, fully Te-anchored anchor (Curti+2020) returns an equal-or-larger deficit, so the signal is not a KE08 artifact. (2) \emph{Orthogonal selection}: the deficit grows, rather than collapses, in a lensed/continuum subsample whose selection is unrelated to emission-line strength, which the emission-line-selection hypothesis cannot reproduce --- \emph{but on $N=4$ galaxies}. (3) \emph{Multi-survey (partial)}: a second reduction (Heintz+23) reports the same-sign offset, \emph{but its only Nakajima-independent objects are those same $4$ lensed galaxies}. The prior companion forward-model, which bounded emission-line selection to $\Delta_{\rm sel}\approx0.10$~dex, now has a model-independent empirical counterpart pointing the same way. Together these make the deficit \emph{selection-robust} and \emph{calibration-robust}, which is a genuine, publishable step beyond a selection-bounded upper limit --- and short of a validated detection by exactly one thing: an independent orthogonal sample larger than $N=4$.

\section{Conclusion}
Once SDSS and the Nakajima+23 $z>7$ census are placed on a common scale and matched in stellar mass, a gas-phase O/H deficit of $\Delta=0.45$~dex (KE08) / $0.55$~dex (Curti+2020) survives on two independent calibration transfers with bootstrap CIs excluding zero ($N=16$; caveat: single emission-line survey). The deficit does not collapse but \emph{grows} to $\Delta=0.68$/$0.78$~dex in an orthogonally-selected $N=4$ lensed subsample, breaking the selection$\leftrightarrow$evolution degeneracy --- but on only four intrinsically metal-poor lensed dwarfs that also constitute the only cleanly-independent second survey, so the independent-survey and orthogonal axes are the same galaxies. The honest statement is therefore: \emph{a selection-robust, multi-calibration $z>7$ mass--metallicity deficit of $\approx0.45$--$0.55$~dex, growing to $\approx0.68$--$0.78$~dex under orthogonal selection, is confirmed on a single small ($N=4$) orthogonal anchor --- strong evidence for genuine early chemical evolution, but not a fully independent detection.} The one decisive next step is a second, larger ($N\gtrsim15$) orthogonally-selected sample from fields not in Nakajima+23; that single addition, and nothing else, converts this into a validated detection.

\section{Caveats}
This is an automated result. (1) The second-survey axis (S) and the orthogonal axis (O) rest on the \emph{same} $4$ lensed galaxies: Heintz's CEERS objects duplicate Nakajima's (redshift-matched $<0.01$), and Curti+24 is unreachable per-object, so the emission-line side is effectively single-survey. (2) $N=4$ is fragile: the bootstrap of four values cannot probe tail behavior, and the CIs and LOO floors exclude zero on four intrinsically metal-poor ($\mathrm{O/H}=7.29$--$7.97$) lensed dwarfs. (3) The direct-Te subset ($N=4$) barely excludes zero and is the more selection-biased channel, so it corroborates but does not independently confirm; the full/CEERS Heintz runs are non-independent sensitivity, not confirmation. (4) SDSS metallicities are fibre-central and reconciled by published-relation transfer, not re-derived from fluxes; no $z>7$ simulation comparator on the matched scale is included. The claim ceiling is a selection-robust, multi-calibration deficit confirmed on a single small orthogonal anchor --- never a validated $z>7$ MZR detection until the independent $N\gtrsim15$ orthogonal sample is in hand.

\begin{acknowledgments}
Based on the public SDSS MPA--JHU catalog, the Nakajima et al.\ (2023) VizieR table \texttt{J/ApJS/269/33}, and the Heintz et al.\ (2023) arXiv:2212.02890 per-object table. This manuscript was generated end-to-end by the autonomous NebulaMind overnight research pipeline as a pre-registered worked example; results are automated, descriptive, and not peer-reviewed.
\end{acknowledgments}

\begin{thebibliography}{}
\bibitem[Curti et al.(2020)]{Curti2020} Curti, M., Mannucci, F., Cresci, G., \& Maiolino, R. 2020, MNRAS, 491, 944
\bibitem[Curti et al.(2024)]{Curti2024} Curti, M., D'Eugenio, F., Carniani, S., et al. 2024, A\&A, 684, A75
\bibitem[Heintz et al.(2023)]{Heintz2023} Heintz, K.~E., Brammer, G.~B., Giménez-Arteaga, C., et al. 2023, Nature Astronomy, 7, 1517
\bibitem[Hirschmann et al.(2023)]{Hirschmann2023} Hirschmann, M., Charlot, S., Feltre, A., et al. 2023, MNRAS, 526, 3610
\bibitem[Kewley \& Ellison(2008)]{KewleyEllison2008} Kewley, L.~J., \& Ellison, S.~L. 2008, ApJ, 681, 1183
\bibitem[Langeroodi et al.(2023)]{Langeroodi2023} Langeroodi, D., Hjorth, J., Chen, W., et al. 2023, ApJ, 957, 39
\bibitem[Nakajima et al.(2023)]{Nakajima2023} Nakajima, K., Ouchi, M., Isobe, Y., et al. 2023, ApJS, 269, 33
\bibitem[Sanders et al.(2021)]{Sanders2021} Sanders, R.~L., Shapley, A.~E., Jones, T., et al. 2021, ApJ, 914, 19
\bibitem[Sanders et al.(2024)]{Sanders2024} Sanders, R.~L., Shapley, A.~E., Topping, M.~W., Reddy, N.~A., \& Brammer, G.~B. 2024, ApJ, 962, 24
\bibitem[Tremonti et al.(2004)]{Tremonti2004} Tremonti, C.~A., Heckman, T.~M., Kauffmann, G., et al. 2004, ApJ, 613, 898
\end{thebibliography}

\end{document}
